\begin{proposition}[Topology of the Weierstrass product]\hypertarget{tag-01BM}{}\label{cdp:st:B3}
\StatementRecord{cdp:st:B3}{proved}{01BM}{\SeriesRef{cdp:st:B2}; \cite{CDP20}}

The threefold $Y$ of Statement \ref{cdp:st:B2} has $H^1(Y,\Z)=0$, $H^2(Y,\Z)\simeq\Z^{11}$, algebraic dimension $3$, and $e(Y)=24$. Its map $g$ is holomorphic, proper, surjective, and has connected fibres. For the line $A$ and the nodal fibre $S_i$ in Statement \ref{cdp:st:B2}, the two conormal-sequence endpoint groups vanish and $h^0(S_i,\Omega_Y^1|_{S_i}\otimes A)=1$. These are the groups used in the fibrewise inference of \cite[Section~7]{CDP20}. The line $M_E$ restricts to a nontorsion line on each nodal cubic. Theorem~2.2 of that paper assumes $H^2(Y,\Z)=0$ and algebraic dimension one.
\end{proposition}


\begin{proof}
\begin{enumerate}[label={$\langle1\rangle$\arabic*.},leftmargin=*]
\item 
The residue calculation above gives $K_E=h^*\OO(-1)$. Relative duality gives $R^1h_*\OO_E=\OO(-1)$, since every fibre is an integral genus-one curve, and $h_*\OO_E=\OO$. Thus $q(E)=p_g(E)=0$ and $H^0(E,K_E^{\otimes m})=0$ for every positive $m$. Castelnuovo's rationality criterion, in the form $q=P_2=0$ for a smooth projective surface, makes $E$ rational. A rational smooth projective surface is obtained from a minimal rational surface by point blowups; its first integral homology vanishes and its second integral cohomology is free. Euler additivity gives $e(E)=12$, since the twelve nodal fibres each have Euler characteristic one and the smooth elliptic fibre has Euler characteristic zero. Poincar\'e duality gives $b_2(E)=10$. The integral K\"unneth isomorphism for $Y=E\times\PP^1$ gives the asserted groups and Euler number. Projectivity and dimension give $a(Y)=3$. Connectedness and surjectivity follow directly from the integral cubic fibres of $h$ and the second product factor.

\item At a nodal fibre write its cubic as $y^2=(x-r)^2(x+2r)$, where $r\ne0$. Its normalization is $x=v^2-2r$, $y=v(v^2-3r)$, and the two preimages of the node are $v=\pm\sqrt{3r}$. At $P$ one has $v_P^2=2r$, while $O$ is the point at infinity. The rational function with divisor $P-O$ on the normalization gives the gluing multiplier $\lambda=(v_P-\sqrt{3r})/(v_P+\sqrt{3r})$, up to inversion. Writing $q=v_P/\sqrt{3r}$ gives $q^2=2/3$ and $\lambda^2+10\lambda+1=0$. Both roots are real and neither is $1$ or $-1$, so $\lambda$ has infinite order. Thus $M_E$ restricts to a nontorsion line on every nodal fibre. The two endpoint vanishings and the one-dimensional ambient cotangent group follow from Statement \ref{cdp:st:B2} and its Serre-duality identification.

\end{enumerate}
\end{proof}